\documentclass[12pt]{article}

\usepackage[margin=1in]{geometry}
\usepackage[T1]{fontenc}
\usepackage[utf8]{inputenc}
\usepackage{microtype}
\usepackage{parskip}
\usepackage{setspace}
\usepackage{titlesec}
\usepackage{hyperref}
\hypersetup{
    colorlinks=false,
    pdfborder={0 0 0},
    pdftitle={Engineering Challenges in Building a Fault-Tolerant Processor},
}

\setstretch{1.15}

\titleformat{\section}{\large\bfseries}{}{0em}{}[\titlerule]
\titlespacing{\section}{0pt}{1.5em}{0.75em}

\title{\textbf{Engineering Challenges in Building a Fault-Tolerant,
Security-Hardened Linux-Capable RISC-V Processor on a 1\,mm$^2$ Die}}
\author{}
\date{}

\begin{document}

\maketitle
\thispagestyle{empty}

% ─────────────────────────────────────────────────────────────
\section*{Feature Categories}

We organize the features we are implementing into seven categories,
which we reference throughout the rest of this document rather than
individual implementations.

\begin{description}
  \setlength{\itemsep}{0.4em}

  \item[\textbf{Error Correction}]
    covers all ECC and parity code mechanisms across cache SRAMs, the
    register file, CSR registers, pipeline registers, and the memory bus.

  \item[\textbf{Redundancy}]
    covers physical and logical duplication of critical state and logic,
    including TMR-style protection of control signals and registers, with
    spatial separation between instances.

  \item[\textbf{Re-execution}]
    covers fault detection through duplicated execution and comparison,
    including a shadow execution pipeline that also serves as our
    execution unit checker, a run-behind hardware thread that verifies
    load and store consistency, and periodic insertion of known-result
    instructions.

  \item[\textbf{Architectural State Hardening}]
    covers protection of privilege-critical registers against physical
    tampering, targeting the CSR file, interrupt controller state, and
    privilege mode registers.

  \item[\textbf{Environmental and Power Hardening}]
    covers circuits that detect and respond to out-of-bounds supply
    conditions.

  \item[\textbf{Security Controls}]
    covers hardware entropy generation and debug interface locking.

  \item[\textbf{Validation Infrastructure}]
    covers the synthesis, simulation, FPGA boot, co-simulation, and
    fuzzing tooling needed to verify all of the above.
\end{description}

\noindent\rule{\linewidth}{0.4pt}

% ─────────────────────────────────────────────────────────────
\section{The Area Constraint}

At 65\,nm TSMC with 1\,mm$^2$ of die area, we are not working with
much margin. A Linux-capable CVW configuration (RV64GC, Sv39 MMU, L1
I/D caches, CLINT, PLIC, FPU) is already a large design; a pair of
16\,KB L1 caches already takes up 150,000 to 200,000\,$\mu$m$^2$ in
bit cell area alone, before any periphery. By the time we have a
functional Linux core, adding the aforementioned features means cutting
or shrinking something else.

Error correction compounds across every structure we apply it to;
re-execution is the most expensive category since it duplicates
execution logic and maintains additional buffering infrastructure;
redundancy adds overhead wherever it is applied. None of these costs are
fixed at design time, and decisions get revisited every synthesis
iteration: which structures get full error correction versus parity, how
aggressively we apply redundancy, whether re-execution covers a full
pipeline instance or a reduced subset. Those runs take hours at this
scale, meaning the synthesis flow has to be treated as actual
infrastructure: automated, run regularly, and tracked across every
feature addition.

% ─────────────────────────────────────────────────────────────
\section{Physical Design}

Even if the total gate count fits, getting everything placed and routed
correctly is its own problem. Redundant instances need spatial
separation; a localized EM pulse or power glitch that hits both copies
simultaneously defeats the redundancy, and the PnR tool has no native
understanding of this. It has to be expressed as explicit placement
constraints or floorplan blockages, alongside timing and routing density
requirements. Error correction check logic ideally sits co-located with
the structures it protects; pushing it far from cache arrays adds
routing overhead directly onto timing-critical read paths. On a
1\,mm$^2$ die, satisfying all of these simultaneously has no guaranteed
clean solution.

% ─────────────────────────────────────────────────────────────
\section{The I/O Ring Bottleneck}

We have at most 53 I/O rings total. After allocating for JTAG (5
signals), power and ground (8 to 12 rings), clock and reset (2 to 3),
and other necessary I/O, we are left with around 30 to 35 rings for
memory. A standard RISC-V cacheline is 512\,bits; that is not enough
to transfer one in a single cycle.

We are still working through what the right memory interface looks like;
the general direction is a serialized burst protocol with a narrow data
bus. Every cache miss pays a multiplied fill penalty as a result, and
the bus width is a hard ceiling on memory performance regardless of how
fast the core runs. If the FPGA build uses a wider memory interface for
convenience, the cache fill timing will be fundamentally different from
what silicon sees, and the FPGA run loses validity as a proxy. While
optimal performance is not the primary focus of this tape-out,
trade-offs between performance and complexity need to be deeply evaluated
for our results to be interesting.

% ─────────────────────────────────────────────────────────────
\section{Verifying Linux Correctness}

Linux boot is somewhere between one and ten billion instructions. A full
RTL simulation in Verilator already takes hours per boot run without
fault injection or co-simulation overhead; the practical signal for
progress is UART output, which tells us the system is mostly working
but gives no instruction-level visibility into whether error correction
fired correctly, whether re-execution caught an injected fault, or
whether architectural state hardening triggered when it should have.

This pushes fault injection and recovery validation onto the FPGA, which
can boot Linux in tens of seconds to minutes at realistic frequencies.
But making the FPGA a meaningful validation artifact means building fault
injection RTL into the build itself, and that fault injection RTL needs
its own verification. Our security controls create an additional
lifecycle problem: once debug locking is engaged, JTAG visibility is
gone; the validation flow has to be structured so that access control
hardening is tested last, or the design needs a secure re-enable
mechanism hardened against the same physical threat model.

% ─────────────────────────────────────────────────────────────
\section{Timing Closure}

Error correction decode logic sits on the load-use latency path;
architectural state hardening checks sit in the execute stage;
re-execution introduces forwarding and buffering paths the base CVW
pipeline was not designed for. At 65\,nm the margins are not as thin as
they would be at an advanced node, but the bigger risk is missing timing
narrowly because a late-added feature closes off the last bit of slack
on a path that was already close. Critical paths need to be tracked
across every synthesis iteration for the same reason area does; we want
to tape out at a reasonable clock frequency even if performance is not
the primary objective.

% ─────────────────────────────────────────────────────────────
\section{Cutting CVW Down to Fit}

CVW was built as a highly parameterizable RV64GC core; that
parameterizability is a liability here. The configuration space is
large: cache sizes, MMU modes, privilege support levels, FPU, debug
module, peripheral controllers. CVW's parameters are deeply
interconnected; changing cache geometry affects MMU page table walks,
modifying peripheral support changes what device drivers Linux can use,
and some parameters produce RTL that simulates fine for simple tests but
breaks under a real Linux workload. The only reliable way to know
whether a configuration reduction is safe is to boot Linux with it,
which brings back the hours-long simulation problem.

Several of our feature categories also require modifying CVW's internal
structures, not just wrapping around them. Architectural state hardening
requires understanding exactly how CVW implements privilege-critical
register reads and writes before we can safely add protection logic.
Re-execution mechanisms have to tap into the CVW pipeline at specific
points that were not documented for external modification. Error
correction on register file reads has to integrate with CVW's forwarding
network, which was not designed to carry check bits alongside data. A
wrong assumption in any of these produces a bug that only surfaces
during Linux boot, hours into simulation.

% ─────────────────────────────────────────────────────────────
\section{The Mechanisms Do Not Compose Cleanly}

The clearest case involves error correction and re-execution operating on
the same data path. If error correction corrects a fault before a value
reaches the execution stage, both pipelines receive the corrected value;
they agree, and the correction happens without raising any flags. That
is correct behavior for a soft error, but it means our re-execution
mechanisms have zero visibility into faults that error correction
handles. Separately, the re-executing path runs behind the main pipeline
and issues memory accesses to verify what it observed; if a write-back
or eviction happens in that window, the re-executing path sees updated
data, producing a spurious divergence that does not indicate any real
fault.

Architectural state hardening protects privilege-critical register
values, but a glitch that corrupts privilege mode state does not
necessarily corrupt an execution unit result; it changes what
instructions are architecturally permitted. Our re-execution mechanism
checks execution unit outputs and does not monitor privilege state,
meaning a successful glitch that forces a privilege transition without
going through a normal register write path would not be caught. The
fault accounting state written when our execution unit checker detects a
failing unit also has to be included in our architectural state hardening
scheme; the fault detection and state hardening infrastructures are not
independent and have to be co-designed, including the full recovery
hierarchy from silent correction through trap or reset, which also has
to interact correctly with the OS.

% ─────────────────────────────────────────────────────────────
\section{The Validation Infrastructure}

Our re-execution mechanism that inserts known-result instructions at
runtime requires that the co-simulation oracle knows exactly which
instructions were inserted and when. The entropy source driving the
insertion has to be observable in simulation; if it is not, co-simulation
breaks at every point an inserted instruction appears, meaning the
insertion hardware has to be co-designed with the validation approach
from the beginning.

Fuzzing for a Linux-capable processor with all of these fault tolerance
categories requires stimulus that targets privilege boundaries, fault
injection timing, and recovery scenarios; standard coverage metrics do
not capture any of that, and off-the-shelf fuzzing tools are not
sufficient. The synthesis flow for area and timing tracking also has to
be automated from the start; if it is run ad hoc, overruns get
discovered late, and getting to tapeout on a 1\,mm$^2$ die requires
continuous visibility into real post-synthesis numbers across every RTL
iteration.

% ─────────────────────────────────────────────────────────────
\section{Development Pipeline as a Response to Simulation Constraints}

A single Linux boot in Verilator takes hours, meaning we cannot block
feature integration on waiting for full simulation runs to complete.
Our response is to pipeline the work: the immediate priority is reducing
the CVW configuration to the minimum needed for Linux boot, freeing up
as much of the 1\,mm$^2$ budget as possible before integration begins.
That work is happening across the team in parallel so we enter the
integration phase with a known baseline rather than discovering overruns
mid-integration.

Once we have that baseline, we run three tracks concurrently: writing
and integrating new features into the RTL, running synthesis on whatever
has been integrated for area estimates, and running simulation on
whatever version of the design has been stable long enough to simulate.
Synthesis gives us area feedback on features before simulation has
validated them; that is an acceptable tradeoff given the constraints.
What this requires is that synthesis results are treated as the primary
feedback signal for area decisions, and that simulation is structured to
catch integration bugs as features stabilize rather than as a gate
before any feature work proceeds.

\end{document}
